\chapter{Implementazione}
\label{chap:implementazione}

\iffalse
Questo capitolo descrive come le scelte del Capitolo~\ref{chap:progettazione} sono state
realizzate, seguendone la stessa articolazione. Vengono riportati i soli
punti in cui la realizzazione ha richiesto decisioni non deducibili dalla progettazione,
ovvero quelli in cui il formato dei dati, il protocollo dei servizi remoti o il meccanismo
di caricamento di Alchemist hanno imposto una forma particolare.
\fi

\section{La lettura dei dati}
\label{sec:rappresentazione-impl}

La lettura di \ac{NetCDF} e \ac{GRIB} è delegata alla libreria NetCDF-Java~\cite{netcdfjava},
sviluppata da Unidata. La libreria espone entrambi i formati attraverso lo stesso modello
di dati, il \ac{CDM}, che descrive un file come un insieme di variabili definite su assi
tipizzati (\req{F8}).

I file sono aperti in modalità \emph{enhanced}, nella quale la libreria interpreta i metadati
già presenti nel file per esporre gli assi tipizzati e, soprattutto, sostituisce
automaticamente i valori convenzionali di riempimento con \texttt{Double.NaN}. La
convenzione interna del modulo per i valori mancanti coincide quindi con quella che la
libreria produce.

\iffalse
\subsection{L'accesso ai formati}
\label{subsec:formati-impl}

La lettura di \ac{NetCDF} e \ac{GRIB} è delegata alla libreria NetCDF-Java~\cite{netcdfjava},
sviluppata da Unidata. La libreria espone entrambi i formati attraverso lo stesso modello
di dati, il \ac{CDM}, che descrive un file come un insieme di variabili definite su assi
tipizzati. Questa uniformità rende possibile \req{F8}: il codice del modulo non distingue
fra i due formati.

I file sono aperti in modalità \emph{enhanced}, nella quale la libreria interpreta i metadati
già presenti nel file per esporre gli assi tipizzati e, soprattutto, sostituisce
automaticamente i valori convenzionali di riempimento con \texttt{Double.NaN}. La
convenzione interna del modulo per i valori mancanti coincide quindi con quella che la
libreria produce.

L'apertura dei file \ac{GRIB} ha richiesto un accorgimento aggiuntivo. Questo formato
necessita di tabelle \ac{XML} esterne
per identificare la variabile di riferimento. La libreria NetCDF-Java contiene
già queste tabelle, ma richiede una specifica realizzazione della fabbrica di parser
per poterle leggere, e fallisce se un'altra risulta registrata a livello di sistema.
L'apertura di un file imposta pertanto temporaneamente la proprietà di sistema
corrispondente, ripristinando il valore precedente, o rimuovendolo se assente,
non appena il file è aperto (\Cref{lst:sax-parser}).

\lstinputlisting[
    language=Kotlin,
    label={lst:sax-parser},
    caption={[Gestione del parser per i file GRIB]Apertura di un file con ripristino della fabbrica di parser preesistente.}
]{listings/implementazione/sax-parser.kt}
\fi

\subsection{L'individuazione e la normalizzazione degli assi}
\label{subsec:normalizzazione-impl}

Il \ac{CDM} associa a ogni asse un tipo che ne specifica il ruolo. La lettura si fonda su
questi tipi anziché sui nomi delle dimensioni, che variano da dataset a dataset: lo stesso
asse compare come \texttt{latitude} in alcuni file e come \texttt{lat} in altri, mentre il
tipo dichiarato è in entrambi i casi quello di latitudine.

Sull'asse temporale il tipo dichiarato non è sempre lo stesso: i dataset che distribuiscono
previsioni espongono l'istante di emissione e quello di validità con tipi distinti, e alcuni
archivi etichettano il proprio unico asse temporale come istante di emissione.
Il \Cref{lst:lettura-assi} riporta l'individuazione dei tre assi, ciascuno ottenuto con
\texttt{findCoordinateAxis} a partire dal tipo cercato.

\lstinputlisting[
    language=Kotlin,
    label={lst:lettura-assi},
    caption={[Individuazione degli assi spaziali e temporali]Individuazione degli assi a partire dal tipo dichiarato nel CDM.}
]{listings/implementazione/readFileAxesAssi.kt}

Per il tempo, l'operatore \texttt{?:} realizza il ripiego appena descritto: si cerca
l'asse di tipo \texttt{Time} e, solo se assente, uno di tipo \texttt{RunTime}. Ogni
ricerca è racchiusa in \texttt{requireNotNull}, così che un asse mancante interrompa
la lettura con un messaggio che nomina il file.
Gli assi vengono richiesti come \texttt{CoordinateAxis1D}, controllando così che siano
monodimensionali.

L'asse temporale viene successivamente ricostruito con \texttt{CoordinateAxis1DTime.factory},
il quale applica le convenzioni \ac{CF}, interpretandone l'unità e l'origine dichiarate nei
metadati. Ciò consente di ottenere istanti assoluti al posto di numeri relativi all'origine.
Eventuali problemi incontrati nella ricostruzione sono raccolti in \texttt{errMsg} e
riportati all'utente.

Lo schema di un file non è però fissato a priori, in due sensi. Il primo è l'ordine delle
dimensioni nella dichiarazione della variabile, che determina la disposizione dei valori in
memoria: le convenzioni \ac{CF} raccomandano
$(\text{tempo}, \text{latitudine}, \text{longitudine})$ senza imporlo~\cite{cf_conventions},
e la libreria restituisce i valori nell'ordine in cui il file li dispone. Poiché le griglie
assumono che le righe corrispondano alle latitudini, un ordine diverso produrrebbe, senza
alcun errore, una griglia con gli assi spaziali scambiati e quindi valori errati (\req{NF5}).
 
Il secondo è la direzione degli assi, che possono essere memorizzati in ordine crescente o
decrescente, mentre l'individuazione di una coordinata è una ricerca binaria
(Sezione~\ref{subsec:ricerca-impl}) che presuppone l'ordine crescente.
 
Il \Cref{lst:permutazione} riporta le due normalizzazioni, che avvengono una sola volta, alla
costruzione della serie, e consentono al resto del modulo di ignorare lo schema originale.

\lstinputlisting[
    language=Kotlin,
    label={lst:permutazione},
    caption={[Lettura di una variabile da un file geospaziale]Lettura di una porzione e suo appiattimento in ordine crescente.}
]{listings/implementazione/gridReading.kt}

In \texttt{readPermutedSlice} la porzione relativa all'istante \texttt{t} viene letta
rispettando l'ordine del file: i vettori \texttt{origin} e \texttt{shape} sono riempiti
usando come indici le posizioni \texttt{timePosition}, \texttt{latPosition} e
\texttt{lonPosition} che le tre dimensioni occupano nella dichiarazione della variabile.
La chiamata a \texttt{permute} riordina poi le dimensioni in
$(\text{tempo}, \text{latitudine}, \text{longitudine})$. Poiché \texttt{permute} restituisce
una vista, la struttura viene copiata tramite \texttt{copy}. Senza la copia, l'accesso per
indice della funzione successiva leggerebbe i valori nell'ordine originale.
 
La funzione \texttt{flattenAscending} appiattisce infine la porzione in un \texttt{DoubleArray} ordinato
per righe di latitudine. Per ogni posizione \texttt{idx} del vettore di destinazione, la
divisione intera e il resto per \texttt{nLon} ricavano gli indici di riga e di colonna; se
un asse è decrescente, l'indice corrispondente viene rovesciato, ad esempio in
\texttt{nLat - 1 - iLat}, prima di leggere il valore sorgente. Il risultato è una griglia
con entrambi gli assi crescenti, indipendentemente da come il file li memorizzava.

\subsection{La selezione della variabile}
\label{subsec:variabile-impl}

Un archivio può contenere più variabili distinte, e la costruzione di una serie richiede di
stabilire quale estrarre (Sezione~\ref{subsec:grid-snapshots}).
Il nome da indicare non è
quello con cui la variabile è catalogata dal servizio, usato durante la selezione dei
parametri in fase di richiesta, ma quello con cui la variabile è identificata all'interno
del file. Quando l'archivio contiene una sola variabile definita sulle tre dimensioni di
interesse, indicarne il nome sarebbe ridondante, e il parametro è per questo facoltativo.
Il \Cref{lst:resolve-variable} riporta i due criteri.

\lstinputlisting[
    language=Kotlin,
    label={lst:resolve-variable},
    caption={Selezione della variabile da estrarre.}
]{listings/implementazione/resolveVariable.kt}

Se il nome è fornito, la variabile viene cercata direttamente con \texttt{findVariable}.
Altrimenti i nomi delle tre dimensioni, ricevuti come argomenti, formano l'insieme
\texttt{targetDims}, e fra le variabili del file vengono selezionate quelle con esattamente
tre dimensioni i cui nomi coincidono con quelle fornite.

I due \texttt{require} permettono di avvertire il modellista se non viene trovata
nessuna variabile specificata o se la risoluzione è ambigua, se più variabili rispettano
tutti i criteri (\req{NF5}).

Se viene correttamente identificata una sola variabile, \texttt{single} restituisce quest'ultima.

\subsection{La costruzione della serie}
\label{subsec:serie-impl}

La realizzazione della serie legge tutti i file di una cartella e li fonde in un'unica
successione. I file vengono elencati in ordine lessicografico e da essi sono esclusi gli
indici che la libreria di lettura genera accanto ai \ac{GRIB}, che risiedono nella stessa
cartella ma non sono dati.

Nella costruzione di \texttt{EagerGridSnapshots} le griglie sono memorizzate in una mappa
ordinata dagli istanti del mondo reale, il che concede due proprietà: l'ordinamento della
successione e l'assenza di istanti duplicati. La fusione rende possibile che due file
coprano il medesimo istante; in tal caso viene conservata la griglia del primo file in
ordine lessicografico e le successive scartate, segnalandolo a chi conduce la simulazione.
L'alternativa, ovvero il fallimento dovuto alla duplicazione degli istanti, sarebbe stata
eccessiva rispetto a una condizione che le richieste ai data store producono spesso.

Ogni griglia letta viene inoltre memorizzata nella realizzazione più adatta alla propria
densità: se questa è inferiore
alla soglia \texttt{SPARSE\_DENSITY\_THRESHOLD}, pari al 10\%, la griglia è costruita come
\texttt{MapRasterGrid}, altrimenti come \texttt{ArrayRasterGrid}.

L'omogeneità dei file è verificata durante la lettura, come mostra il \Cref{lst:eager-init}.

\lstinputlisting[
    language=Kotlin,
    label={lst:eager-init},
    caption={Costruzione della serie dai file di una cartella.}
]{listings/implementazione/EagerGridInit.kt}

La variabile \texttt{reference}, che specifica le informazioni sulla prima griglia
della serie, è inizialmente nulla. Al primo file, l'operatore \texttt{?:}
la inizializza con gli assi appena letti; per ogni file successivo, \texttt{?.also} invoca
invece \texttt{requireMatches}, che confronta assi spaziali e nome della variabile con quelli
di riferimento. Un disallineamento produce così un errore al caricamento della simulazione.
Ogni file è aperto con \texttt{use}, che ne garantisce la chiusura anche in caso di errore,
e le griglie lette da \texttt{readTimestepsFrom} confluiscono nella \texttt{TreeMap}.
Alla fine delle operazioni di lettura, chiavi e valori della mappa diventano gli
istanti e le griglie della serie, già ordinati.

\section{Le strategie di risoluzione}
\label{sec:strategie-impl}

\subsection{La ricerca sugli assi}
\label{subsec:ricerca-impl}

Le tre operazioni di cui le strategie e il layer hanno bisogno sono raccolte in un unico
file di funzioni interne: l'individuazione dell'indice più vicino a una coordinata, quella
della coppia di indici che la racchiude, e il calcolo del peso che ne esprime la posizione
relativa fra i due. Tutte estrapolano in modo costante oltre gli estremi dell'asse,
restituendo l'indice estremo più vicino: è questa proprietà a realizzare sull'asse
temporale l'estrapolazione discussa nella Sezione~\ref{sec:corrispondenza-temporale}.
Il \Cref{lst:axes-search} riporta la prima delle tre.

\lstinputlisting[
    language=Kotlin,
    label={lst:axes-search},
    caption={Individuazione dell'indice più vicino a una coordinata sull'asse.}
]{listings/implementazione/Axes.kt}

La funzione si appoggia a \texttt{Arrays.binarySearch}, che restituisce l'indice della
coordinata se questa è presente sull'asse, altrimenti il valore $-p-1$, dove $p$ è il
punto di inserimento, ovvero l'indice del primo elemento maggiore della coordinata.
\texttt{upperIndex} ricava $p$ da questa codifica, e i rami del \texttt{when} coprono i
quattro casi possibili:

\begin{itemize}
    \item la coordinata è presente sull'asse, e il suo indice viene restituito direttamente;
    \item la coordinata precede il primo elemento, e viene restituito l'indice
    iniziale;
    \item la coordinata supera l'ultimo elemento, e viene restituito l'ultimo indice;
    \item la coordinata si trova tra due elementi dell'asse, e fra i due indici adiacenti viene scelto quello
    più vicino.
\end{itemize}

Il secondo e il terzo ramo sono quelli che realizzano l'estrapolazione costante.

Nessuna di queste funzioni verifica che l'asse ricevuto sia strettamente crescente. 
Il controllo non viene effettuato poiché sarebbe ridondante visto che la monotonia
degli assi spaziali è garantita dalla costruzione delle griglie; la monotonia dell'asse
temporale è garantita invece dalla costruzione della serie. Inoltre,
tale controllo porterebbe a costo lineare un'operazione altrimenti logaritmica.

\subsection{Uguaglianza fra le strategie}
\label{subsec:composizione-impl}

Le strategie prive di configurazione, come \texttt{BilinearInterpolation} o
\texttt{LinearInterpolation}, ridefiniscono l'uguaglianza per soddisfare un vincolo del
meccanismo di caricamento. Quando un parametro è a sua volta un oggetto descritto nel file
YAML, il caricatore deve sapere di quale tipo debba essere l'oggetto, e questo dipende da
quale costruttore del layer verrà invocato. Il caricatore costruisce allora il parametro
una volta per ciascun costruttore possibile e confronta le istanze ottenute: se non
risultano tutte uguali, considera la risoluzione ambigua e interrompe il caricamento.
 
Tutti e tre i costruttori del layer dichiarano la coppia di strategie per l'interpolazione
e la conversione del valore (Sezione~\ref{subsec:costruttori-impl}), per cui una stessa
strategia viene costruita più volte. Il confronto predefinito fra oggetti guarda però
l'identità dei riferimenti, e istanze costruite separatamente non risultano mai uguali fra
loro: il caricamento fallirebbe. Per questo tali strategie estendono una classe base comune,
\texttt{StatelessStrategy}, riportata nel \Cref{lst:stateless-strategy}.

\lstinputlisting[
    language=Kotlin,
    label={lst:stateless-strategy},
    caption={Criterio di identità comune alle strategie prive di configurazione.}
]{listings/implementazione/StatelessStrategy.kt}

La classe ridefinisce \texttt{equals} e \texttt{hashCode} in modo che l'uguaglianza
tra due strategie senza stato sia data solamente dalla classe che istanziano, confrontando
con \texttt{other::class}.

Il \Cref{lst:trilinear} mostra la classe in uso in \texttt{TrilinearInterpolation}.

\lstinputlisting[
    language=Kotlin,
    label={lst:trilinear},
    caption={L'interpolazione trilineare per delegazione.}
]{listings/implementazione/TrilinearInterpolation.kt}

La classe \texttt{TrilinearInterpolation} estende \texttt{StatelessStrategy} e implementa
\texttt{SpatioTemporalInterpolation} delegando, tramite \texttt{by}, a una
\texttt{SeparableSpatioTemporalInterpolation} configurata con
\texttt{BilinearInterpolation} e \texttt{LinearInterpolation}. L'uguaglianza
ereditata resta corretta anche qui, perché la configurazione è stabilita dalla classe stessa
e non può variare fra un'istanza e l'altra.

\section{Il layer}
\label{sec:layer-impl}

\subsection{La corrispondenza temporale}
\label{subsec:corrispondenza-impl}

La corrispondenza fra i due domini temporali è stabilita alla costruzione del layer, come
descritto nella Sezione~\ref{sec:corrispondenza-temporale}. Il \Cref{lst:layer-time} ne
riporta la realizzazione.

\lstinputlisting[
    language=Kotlin,
    label={lst:layer-time},
    caption={Conversione degli istanti della serie in tempo di simulazione.}
]{listings/implementazione/CopernicusLayerTime.kt}

Il blocco \texttt{init} verifica due condizioni: che la serie non sia vuota, poiché una
serie priva di istanti non individua alcun valore, e che \texttt{timeScale} sia una durata
finita e positiva, poiché una scala nulla renderebbe la conversione indefinita e una
negativa invertirebbe l'ordine della successione.
 
La proprietà \texttt{sliceTimes} contiene il risultato della conversione
da istanti reali a tempi di simulazione. L'origine è quella
dichiarata oppure, tramite l'operatore \texttt{?:}, il primo istante della serie, e ogni
istante viene tradotto da \texttt{toSimulationTime}. La funzione si riduce all'espressione
\texttt{(this - origin) / scale}: la differenza fra due \texttt{Instant} produce una
\texttt{Duration}, e il rapporto fra due \texttt{Duration} un numero reale, che è
esattamente il tempo di simulazione cercato.
Il risultato è conservato in un \texttt{DoubleArray}, sul quale opera la ricerca della
Sezione~\ref{subsec:ricerca-impl}.

\subsection{L'interrogazione}
\label{subsec:interrogazione-impl}

L'interrogazione segue la sequenza della \Cref{fig:layer-sequenza}, e il
\Cref{lst:layer-get-value} ne riporta la realizzazione.

\lstinputlisting[
    language=Kotlin,
    label={lst:layer-get-value},
    caption={Risoluzione di un valore a partire da una posizione geografica.}
]{listings/implementazione/CopernicusLayerGetValue.kt}

La prima istruzione di \texttt{getValue} legge il tempo di simulazione corrente. Il tempo
non è però sempre disponibile (Sezione~\ref{subsec:alchemist-time}): la simulazione esiste
solo dopo che il modello è stato costruito, e \texttt{simulationOrNull} restituisce
\texttt{null} se il layer viene interrogato prima. In tal si assume lo zero,
ovvero il valore che la simulazione esporrebbe non appena costruita.
 
\texttt{bracketIndices} restituisce poi gli indici dei due campioni di \texttt{sliceTimes}
che racchiudono \texttt{t}. Quando i due indici coincidono, perché \texttt{t} cade
esattamente su una griglia oppure, per l'estrapolazione costante, al di fuori
dell'intervallo coperto, \texttt{sampleExactSlice} usa due volte la stessa griglia con peso
nullo; altrimenti \texttt{weight} calcola la posizione relativa di \texttt{t} fra i due
campioni.
 
In entrambi i casi si converge in \texttt{sample}, che realizza la parte
finale della sequenza: ottiene le due griglie dalla serie, verifica con \texttt{require} che
la posizione ricada nella regione coperta, segnalando in caso contrario l'errore di
modellazione, e restituisce
\texttt{converter.convert(interpolation.interpolate(\dots))}. In quest'ultima espressione si
nota direttamente la separazione fra la produzione del valore numerico e la sua traduzione
nel tipo esposto. La verifica dei limiti spaziali avviene sulla sola \texttt{gridBefore}, poiché
l'omogeneità degli assi spaziali è accertata alla costruzione della serie.

\subsection{I costruttori e la configurazione}
\label{subsec:costruttori-impl}

Alle tre modalità di costruzione descritte nella Sezione~\ref{subsec:layer-costruzione} corrispondono
altrettanti costruttori di \texttt{DoubleCopernicusLayer}, la specializzazione concreta di
\texttt{CopernicusLayer<T>}, riportati nel \Cref{lst:costr-layer}. La forma dei loro
argomenti non è data però dalle sole modalità di costruzione: in più punti è il meccanismo
di caricamento (Sezione~\ref{subsec:alchemist-loader}) a dettarla.

\lstinputlisting[
    language=Kotlin,
    label={lst:costr-layer},
    caption={[I tre costruttori di DoubleCopernicusLayer.]I tre costruttori di \texttt{DoubleCopernicusLayer}.}
]{listings/implementazione/DoubleCopernicusLayer.kt}

Il primo costruttore non è raggiungibile dalla configurazione, poiché richiede una serie di
griglie già costruita, che il meccanismo di caricamento non è in grado di produrre; non porta infatti
l'annotazione \texttt{@JvmOverloads}, che ne genererebbe le varianti. È la via attraverso la
quale il comportamento del layer è verificabile in isolamento (\req{NF3}).

Il secondo costruttore assume che i file siano già presenti sulla macchina e ne riceve la
collocazione; tutti i file della cartella contribuiscono a formare un'unica serie. I restanti
argomenti ne indicano l'interpretazione: la scala temporale, con valore predefinito di
un'ora reale per unità di tempo simulato; l'origine, che con valore nullo di default indica
il primo istante esposto dalla serie; e il nome della variabile, anch'esso nullo per default
e in tal caso rilevato automaticamente. Tre di questi denotano tipi che il meccanismo di
caricamento non sa costruire, ovvero \texttt{Path}, \texttt{Instant} e \texttt{Duration}, e
sono pertanto dichiarati come stringhe e convertiti all'interno del costruttore.

La durata e l'istante sono convertiti secondo il formato definito dallo standard ISO
8601~\cite{iso8601}, compatto e convertibile senza ambiguità. Ad esempio, un istante può
essere scritto come \texttt{2002-09-30T00:00:00Z} mentre una durata come \texttt{PT6H}. La
conversione nel costruttore ha inoltre l'effetto di rilevare nel caricamento
valori non validi, coerentemente con \req{NF6}.

Il terzo costruttore riceve la descrizione di come ottenere i dati. Quattro argomenti sono
obbligatori: il data store da interpellare, il dataset, la collocazione del file contenente
la selezione dei parametri e l'abilitazione della verifica del digest MD5 sui file scaricati.
Seguono, nell'ordine, gli argomenti del costruttore precedente, poiché una volta ottenuti i
file il comportamento è il medesimo; successivamente i tre destinati al recupero o produzione in cache: la
cartella radice, la collocazione del file delle credenziali, che per default è quella
prevista dal client ufficiale, e una durata limite entro la quale al provider è consentito
provare a popolare la cartella con i dati. Il file delle credenziali viene letto solo nel momento in cui
una richiesta di rete si rende necessaria: una simulazione che trovi i propri dati già in
cache può quindi essere eseguita anche in assenza di credenziali valide.

Due argomenti dell'ultimo costruttore meritano una spiegazione. Il primo è la selezione dei
parametri, che risiede in un file JSON separato anziché nella configurazione YAML. Questo
poiché il meccanismo di caricamento attribuisce a certi nomi un significato proprio, ad
esempio \texttt{type}, con cui è indicata la classe da istanziare, e lo stesso nome può
comparire fra i parametri ammessi da alcuni dataset, dove designa ad esempio il tipo di
rianalisi richiesta. 
\iffalse
Una selezione dichiarata direttamente nella configurazione non è quindi
possibile, poiché il meccanismo interpreterebbe quella chiave come propria ed esigerebbe che
il valore associato nomini una classe.
\fi

Il secondo è l'abilitazione della verifica del digest MD5, mantenuta obbligatoria
benché un valore predefinito sarebbe naturale. Ogni argomento facoltativo genera una
variante della firma, ottenuta rimuovendo gli argomenti a partire dall'ultimo: la
variante più corta di un costruttore è quindi formata dai soli argomenti obbligatori.
Rendendo facoltativo anche questo argomento, la firma più corta del terzo costruttore
sarebbe \texttt{(Environment, String, String, String)}, la stessa che il secondo
produce omettendo \texttt{variable}, \texttt{interpolation} e \texttt{converter}.
Sulla \ac{JVM} le due firme sarebbero
indistinguibili e la classe non compilerebbe. Nemmeno collocare l'argomento altrove
fra i facoltativi risolverebbe, poiché la variante più corta dipende dai soli
argomenti obbligatori. Mantenerlo obbligatorio è dunque l'unico modo di conservare
distinte le due modalità di costruzione.

L'ordine degli argomenti non è indifferente. Di quelli facoltativi può essere omesso solo un
suffisso (Sezione~\ref{subsec:alchemist-loader}), per cui sono dichiarati in ordine
decrescente di probabilità d'uso. Ad esempio, i tre destinati alla cache seguono quelli
sull'interpretazione dei dati perché è assai più frequente dover cambiare la scala temporale
che non la collocazione della cache.

L'environment, dichiarato da ciascuno dei tre, non compare invece fra gli argomenti della
configurazione YAML, essendo iniettato automaticamente.

Il \Cref{lst:yaml} mostra un esempio della forma finale assunta dal layer dal punto di vista di chi
conduce una simulazione (\req{F11}). Si noti che
gli argomenti successivi all'interpolazione sono omessi e assumono quindi i propri
valori di default.

\lstinputlisting[
    language=yaml,
    label={lst:yaml},
    caption={Dichiarazione del layer in un file di simulazione.}
]{listings/implementazione/yaml.yml}

\iffalse
La selezione dei parametri risiede nel file indicato, il cui contenuto è quello atteso dal
servizio ed è riportato nel \Cref{lst:inputs}.

\lstinputlisting[
    language=json,
    label={lst:inputs},
    caption={Selezione dei parametri della richiesta, in un file separato.}
]{listings/implementazione/2mt.json}
\fi

\section{Il sottosistema di acquisizione}
\label{sec:acquisizione-impl}

\subsection{L'identità di una richiesta}
\label{subsec:nome-cartella-impl}

L'identità di una richiesta è
data dal data store, dal dataset e dalla selezione dei parametri; di quest'ultima conta
il contenuto e non l'ordine con cui i parametri compaiono. La selezione è
letta come struttura generica di mappe e liste, e la codifica ne compie una sola
normalizzazione, riportata nel \Cref{lst:canonical-json}.

\lstinputlisting[
    language=Kotlin,
    label={lst:canonical-json},
    caption={Ordinamento della selezione dei parametri.}
]{listings/implementazione/CanonicalJson.kt}

La funzione \texttt{canonicalize} visita ricorsivamente la struttura. Ogni mappa viene
ricopiata in una \texttt{TreeMap}, che mantiene le chiavi ordinate, normalizzandone a sua
volta i valori; le liste sono ricostruite conservandone l'ordine, poiché nelle richieste
ai data store l'ordine di una
lista è semantico, come nel caso di una sequenza di coordinate; ogni altro
valore è restituito invariato. Il metodo \texttt{encode} serializza infine il risultato nella rappresentazione
JSON corrispondente, che per
costruzione è lo stesso per selezioni equivalenti ma con diverso ordine.
 
Il \Cref{lst:coper-req} mostra come \texttt{CopernicusRequest} traduca l'identità in un nome
di cartella.

\lstinputlisting[
    language=Kotlin,
    label={lst:coper-req},
    caption={Derivazione del nome della cartella di cache.}
]{listings/implementazione/CopernicusRequest.kt}

In \texttt{toDirectoryName}, la mappa passata a \texttt{CanonicalJson.encode} contiene le tre
componenti dell'identità, per cui anche \texttt{endpoint} e \texttt{dataset} contribuiscono
all'hash e due richieste uguali verso data store diversi non condividono la cartella. La
codifica è data come argomento alla funzione di hash SHA-256~\cite{sha256}, calcolata da
\texttt{sha256Hex}, e del digest si conservano le prime \texttt{HASH\_PREFIX\_LENGTH}, ovvero
16, cifre esadecimali. Il nome finale della voce in cache è dato dal nome del dataset,
reso sicuro per il file system da \texttt{toFileSystemSafe} che mantiene i soli caratteri ammessi
e sostituisce tutti gli altri con un trattino basso, seguito dalle cifre del digest.
 
Il nome è un compromesso fra leggibilità e probabilità di collisione: la parte derivata dal
dataset consente a chi conduce la simulazione di riconoscere la cartella ispezionando la
cache, mentre le cifre del digest ne stabiliscono l'identità. La lunghezza scelta corrisponde
a 64 bit, quindi due richieste distinte finiscono nella stessa cartella con probabilità
trascurabile.

\subsection{L'atomicità della cache}
\label{subsec:atomicita-impl}

Il requisito \req{NF8} richiede che una voce di cache non sia mai osservabile in stato
parziale. La realizzazione distingue due forme di concorrenza, che affronta con mezzi
diversi: fra thread di una stessa esecuzione e fra esecuzioni distinte del simulatore.

All'interno di una stessa esecuzione la produzione di una voce avviene al più una volta,
secondo la logica riportata nel \Cref{lst:cache-lock}.

\lstinputlisting[
    language=Kotlin,
    label={lst:cache-lock},
    caption={Accesso a una voce di cache da parte di più thread.}
]{listings/implementazione/FileSystemCacheManagerLock.kt}

\texttt{getOrProduce} verifica per prima cosa la presenza della cartella:
nel caso di una voce già prodotta da un'esecuzione
precedente, il costo è un singolo accesso al file system. Se la voce è assente,
\texttt{computeIfAbsent} ottiene il monitor associato alla cartella di destinazione.
L'operazione è atomica sulla \texttt{ConcurrentHashMap}, per cui tutti i thread che cercano
la stessa voce ricevono lo stesso oggetto. La mappa risiede nel \texttt{companion object},
e il monitor è quindi condiviso anche fra istanze diverse di
\texttt{FileSystemCacheManager}.

Dentro il blocco \texttt{synchronized} la presenza della cartella viene verificata una
seconda volta. Il primo thread a entrare non la trova e invoca \texttt{produce}; chi era in
attesa, entrando dopo, la trova prodotta e la restituisce senza scaricare nulla. Se invece la
produzione fallisce, nessuna cartella viene pubblicata, e il thread successivo tenta a sua
volta la produzione. Voci in cache distinte hanno monitor distinti, e
possono quindi essere prodotte in parallelo.

Fra esecuzioni distinte del simulatore non ci possono essere monitor condivisi, e la garanzia
di validità della voce viene ottenuta dal file system: la produzione avviene in una cartella
temporanea, resa visibile sotto il nome definitivo tramite una singola operazione atomica solo dopo che i dati sono stati
correttamente acquisiti. Due esecuzioni concorrenti
possono così sovrapporsi: entrambe scaricano, ma solo la prima a concludere pubblica la
propria cartella, mentre la seconda riconosce che la voce esiste già e scarta il proprio
lavoro. Una voce già presente è considerata completa: non sono previsti recupero da stati
parziali né riverifica del contenuto. Il \Cref{lst:cache-manager} riporta la produzione e la
pubblicazione di una voce.

\lstinputlisting[
    language=Kotlin,
    label={lst:cache-manager},
    caption={Produzione atomica di una voce di cache.}
]{listings/implementazione/FileSystemCacheManager.kt}

\texttt{produce} crea la cartella temporanea con \texttt{Files.createTempDirectory},
la quale popola tramite la fonte con \texttt{provider.fetch}.
Prima della pubblicazione, \texttt{check(hasData(temp))} esclude che
una fonte terminata senza errori ma senza aver prodotto file generi una voce vuota.

Il metodo \texttt{promote} tenta quindi lo spostamento con l'opzione \texttt{ATOMIC\_MOVE}. 
Se fallisce e la cartella di destinazione esiste, allora significa che un'altra
esecuzione ha pubblicato per prima, e restituisce \texttt{false}; se fallisce e la cartella
non esiste, l'errore è di altra natura e l'eccezione viene rilanciata.

Il flag \texttt{promoted} unisce i casi nel blocco \texttt{finally}: la cartella temporanea
viene eliminata sia quando si è verificato un errore, inclusi i fallimenti della fonte, sia
quando la corsa è stata persa, così che nessuna voce incompleta resti nella cache. In
entrambi i casi di successo \texttt{produce} restituisce la stessa cartella definitiva, e
chi la riceve non ha bisogno di sapere chi l'abbia pubblicata.

\subsection{Il ciclo di vita di una richiesta}
\label{subsec:ciclo-richiesta-impl}

Il dialogo con i data store segue il paradigma descritto nella \Cref{sec:ogc-api-processes}.
Il riferimento dei servizi\footnote{ECMWF Data
Stores Processing API: \url{https://cds.climate.copernicus.eu/api/retrieve/v1/docs}.} ne
descrive gli endpoint e lo schema dei messaggi, ma non come questi vengano popolati: diversi
aspetti del comportamento concreto sono stati quindi ricostruiti sperimentalmente, osservando
le risposte e confrontandole con il codice del client ufficiale distribuito da
ECMWF. Il modulo confina l'intero dialogo in
\texttt{CopernicusDataStoreProvider}, che è l'unico punto in cui la struttura del protocollo
è nota.
 
Una richiesta viene sottomessa in \texttt{POST} al percorso di esecuzione del processo, che
lo standard definisce come \texttt{/processes/\{processID\}/execution}. I data store lo
precedono con un prefisso proprio, a sua volta preceduto dall'indirizzo del data store
ricevuto dal layer come \texttt{endpoint}, e assegnano al \emph{processID} il nome del
dataset richiesto. Il corpo della richiesta contiene la sola selezione dei parametri,
associata alla chiave \texttt{inputs}. L'identità dell'utente viaggia nell'intestazione
\texttt{PRIVATE-TOKEN} in tutte le chiamate al servizio, tranne il recupero finale del file,
che avviene senza autenticazione.
 
La risposta alla sottomissione espone l'indirizzo da interrogare in tre forme: un
identificativo del processo, l'intestazione \texttt{Location} e un elenco di collegamenti.
Il client segue il collegamento, sulle orme del client ufficiale. Le fasi successive procedono
allo stesso modo, seguendo il collegamento di monitoraggio durante l'attesa e
quello dei risultati a completamento avvenuto.
 
Lo stato di una richiesta viene letto come stringa e non come valore di un insieme chiuso,
come fa anche il client ufficiale, così che uno stato non previsto produca un avviso
al posto di interrompere l'attesa. Il modulo li ripartisce in tre categorie:
 
\begin{itemize}
    \item \texttt{successful}, unico stato di successo, dopo il quale il collegamento ai
    risultati è disponibile;
    \item \texttt{accepted} e \texttt{running}, stati transitori durante i quali l'attesa
    prosegue;
    \item \texttt{failed}, \texttt{rejected}, \texttt{dismissed} e \texttt{deleted}, stati
    terminali di fallimento, per i quali la causa va ricostruita e mostrata al modellista
    (Sezione~\ref{subsec:fallimenti-impl}).
\end{itemize}
 
L'attesa è realizzata da \texttt{awaitSuccess}, riportata nel \Cref{lst:provider}.
 
\lstinputlisting[
    language=Kotlin,
    label={lst:provider},
    caption={[Attesa del completamento di una richiesta.]Attesa del completamento di una richiesta. Il ramo degli stati transitori,
             omesso, segnala periodicamente che l'attesa prosegue.}
]{listings/implementazione/DataStoreProvider.kt}
 
La scadenza complessiva è calcolata una sola volta tramite \texttt{TimeSource.Monotonic}.
Il limite è di trenta minuti per default
ed è configurabile dal layer (Sezione~\ref{subsec:costruttori-impl}), poiché la durata di
un'elaborazione dipende dall'ampiezza della richiesta.
 
A ogni iterazione, lo stato viene estratto dal corpo della risposta di monitoraggio con
\texttt{parseStatus}, e i rami del \texttt{when} realizzano le tre categorie. Lo stato
\texttt{successful} restituisce l'indirizzo dei risultati, e se il collegamento corrispondente
manca la risposta è incoerente e produce un errore. Gli stati terminali passano
a \texttt{failOnStatus}, che ricostruisce la causa del fallimento
(Sezione~\ref{subsec:fallimenti-impl}). Il ramo \texttt{else} registra infine un avviso per
uno stato sconosciuto, senza interrompere l'attesa.
 
Superato il controllo della scadenza con \texttt{check}, il thread attende e l'intervallo
viene raddoppiato, limitato da \texttt{coerceAtMost} al valore di \texttt{maxPollInterval}:
partendo da due secondi fino a un massimo di due minuti, le interrogazioni diminuiscono per
le elaborazioni lunghe, evitando chiamate di rete inutili. Proprio perché le interrogazioni
diventano sempre più rade, il ramo degli stati transitori segnala a intervalli regolari che
l'attesa prosegue, così che un'esecuzione lunga non appaia bloccata.

\subsection{La verifica di integrità}
\label{subsec:integrita-impl}

Fra le proprietà esposte prima della consegna finale, la dimensione in byte e il digest MD5
consentono di accertare che il file ricevuto sia quello aspettato. La verifica avviene sul
file scaricato. Se il file consegnato è un archivio, il digest e la dimensione si riferiscono
a questo piuttosto che ai file che contiene, quindi la verifica viene effettuata prima
di eventuali azioni di estrazione.

Le due proprietà sono esposte con i nomi \texttt{file:size} e \texttt{file:checksum} definiti
dalla File Info Extension di STAC~\cite{stacfile}, che però specifica per il checksum
un multihash, ossia un digest preceduto dall'identificatore dell'algoritmo e dalla
propria lunghezza. Il valore dato dai data store è invece il digest MD5 grezzo.

Il digest esposto presenta una particolarità non documentata: è rappresentato come numero
esadecimale privo delle cifre nulle iniziali, e ha quindi lunghezza variabile, mentre il
digest calcolato localmente ne ha sempre trentadue. Un confronto diretto fra le due stringhe
fallirebbe quindi ogni volta che il digest del file comincia per zero, pur essendo il file
integro e correttamente leggibile. Il valore ricevuto viene pertanto riportato alla lunghezza
standard prima del confronto.

Anche riportando il digest alla lunghezza standard, in due occasioni nel corso dello sviluppo il confronto è
fallito su file la cui dimensione coincideva con quella dichiarata e che la libreria di
accesso ai formati leggeva senza errori: i dati acquisiti erano dunque utilizzabili. Non
avendo potuto stabilire se la causa risieda nel modo in cui il servizio calcola il digest,
la verifica di quest'ultimo è stata resa disattivabile, così che un'anomalia di
questo tipo non impedisca di eseguire una simulazione con dati validi. La dimensione,
unico controllo che il client ufficiale effettua, resta invece sempre verificata.

\subsection{Le modalità di fallimento}
\label{subsec:fallimenti-impl}

Un tentativo di acquisizione può fallire in due modi distinti, che il servizio segnala in
forme diverse, e la ricostruzione della causa è la parte che si è discostata
maggiormente dal protocollo descritto nella Sezione~\ref{sec:ogc-api-processes}.
 
Il primo è un errore a livello di protocollo, segnalato da un codice di stato di errore e
accompagnato da un corpo conforme all'RFC 7807~\cite{rfc7807}. La convenzione prevede che
la spiegazione leggibile risieda nel campo \texttt{detail}; i data store, però, lo popolano
solo in parte, e collocano la traccia di esecuzione in \texttt{traceback}, un membro di
estensione.
 
Il secondo è il fallimento dell'elaborazione: qui il protocollo ha successo, la risposta ha
esito positivo, e il fallimento è dichiarato dallo stato. Lo standard prevede in questo caso
un campo di primo livello destinato al messaggio, che però i data store lasciano vuoto. La
causa è raggiungibile per un'altra via: il collegamento ai risultati, che compare anche su
una richiesta fallita, risponde con un errore il cui corpo è un problema descritto secondo la
convenzione precedente, ed è lì che la traccia di esecuzione si trova. Il client segue quindi
quel collegamento, e ricade sul campo previsto dallo standard solo se questa via non produce
nulla, nel caso in cui una futura revisione del servizio lo popoli.
 
In entrambi i casi il messaggio mostrato al modellista è prodotto da \texttt{describe},
riportata nel \Cref{lst:problem-detail}.
 
\lstinputlisting[
    language=Kotlin,
    label={lst:problem-detail},
    caption={[Ricostruzione del fallimento di una richiesta]Ricostruzione della causa di un fallimento dal corpo della risposta.}
]{listings/implementazione/ProblemDetail.kt}
 
La prima istruzione realizza l'ordine di precedenza con una catena di operatori \texttt{?:}:
viene usato \texttt{detail} se presente, altrimenti \texttt{traceback}, altrimenti il titolo
del problema. Se nessuno dei tre è presente la funzione restituisce una stringa vuota.
Al messaggio vengono poi accodati il tipo
del problema e, se presente, l'identificativo fornito dal servizio, che permette
di ricondurre l'errore alla specifica richiesta.
